% Source: https://artofproblemsolving.com/community/c3h1216440_1st_reun_mock_middle_school_math_contest
% Stable identity: aops-topic-1216440-q10
The radius is $6$, so the whole circle has area $36\pi$.

\begin{asy}
usepackage("amsmath");
size(3.5cm);
defaultpen(fontsize(10pt));
for(int k=0;k<36;++k) draw((0,0)--6*dir(10*k),gray+0.3);
fill((0,0)--arc((0,0),6,0,30)--cycle,lightblue);
for(int k=0;k<4;++k) draw((0,0)--6*dir(10*k));
draw(circle((0,0),6)); label("$6$",(0,0)--6*dir(200),NW);
label("\mbox{}",(min(currentpicture,user=true)+max(currentpicture,user=true))/2,black); // keeps the text after the figure black when LaTeX typesets the labels inline
\end{asy}

Three of the $36$ equal sectors make up $\dfrac{3}{36}$ of the circle, so their area is $\dfrac{3}{36}\cdot36\pi=\boxed{3\pi}$.
